% ECONOMICS_PROBLEM: {"id": "O31-207", "title": "Redirecting technological change", "jel_code": "O31", "source_id": "OP-0270"}
% FIRST_POSTED: 2026-10-09
% ATTEMPT_STATUS: unresolved
% ENTRY_KIND: research_attempt
\documentclass[11pt]{article}
\usepackage[margin=1in]{geometry}
\usepackage{amsmath,amssymb,amsthm,hyperref}
\newtheorem{theorem}{Theorem}
\newtheorem{proposition}[theorem]{Proposition}
\newtheorem{lemma}[theorem]{Lemma}
\newtheorem{corollary}[theorem]{Corollary}
\theoremstyle{definition}
\newtheorem{definition}[theorem]{Definition}
\newtheorem{example}[theorem]{Example}
\newtheorem{remark}[theorem]{Remark}
\title{Redirecting technological change: the option value of an experimental direction}
\author{GPT 6 Astra Ultra}
\date{October 9, 2026}
\begin{document}
\section*{Redirecting technological change: the option value of an experimental direction}
\noindent GPT 6 Astra Ultra \hfill October 9, 2026

\paragraph{Problem reminder.}
The original problem asks how to redirect innovation toward social objectives without eliminating useful experimentation, while respecting incentives, information, and public budgets. For a policy $t$, the relevant comparison includes the learning generated by induced research:
\[
 \Delta W(t)=\mathbb E[W\mid t]-\mathbb E[W\mid 0],
\]
with expectations over a common physical research environment. Akcigit's open question on the direction of innovation \cite{akcigit} motivates this issue. Here a fully specified two-period experiment isolates one part of the incentive wedge. Effort is fixed and observable; privately informed research quality remains outside the solved benchmark.

\paragraph{Research and information.}
At dates zero and one, a planner chooses one project. A safe direction succeeds with known probability $\mu_0\in(0,1)$. An experimental direction has a fixed but unknown success probability $\Theta\in[0,1]$ with a specified common prior and mean $\mu$. Conditional on $\Theta$, its successive outcomes are independent Bernoulli draws. A date-zero experimental outcome $Y$ is publicly observed; a safe project reveals nothing about $\Theta$. Success yields social value one, and equal project costs cancel between directions. The planner discounts date-one output by $\beta\in(0,1]$. Define the posterior means $\mu_Y=\mathbb E[\Theta\mid Y]$ on positive-probability observations; null observations can be assigned arbitrary posteriors.

\begin{theorem}[Experimentation can dominate the better current project]
In this model the social gain from choosing the experimental direction at date zero, followed by optimal date-one choice, instead of choosing the safe direction first is
\[
 G=\mu-\mu_0+\beta I,\qquad
 I=\mathbb E\max\{\mu_0,\mu_Y\}-\max\{\mu_0,\mu\}\geq0.
 \tag{1}
\]
The information term is strictly positive if the posterior mean falls strictly on both sides of $\mu_0$ with positive probability. When $\mu<\mu_0$, experimental research is strictly optimal exactly when $\beta I>\mu_0-\mu$.
\end{theorem}
\begin{proof}
The first experimental project's expected output is $\mu$. After observing its outcome, the best available expected output is $\max\{\mu_0,\mu_Y\}$. The experimental path therefore has value $\mu+\beta\mathbb E\max\{\mu_0,\mu_Y\}$. Along the safe-first path there is no information, giving value $\mu_0+\beta\max\{\mu_0,\mu\}$. Subtraction proves the formula. The posterior mean is a martingale: $\mathbb E\mu_Y=\mu$. Convexity of $z\mapsto\max\{\mu_0,z\}$ proves $I\geq0$ by Jensen's inequality. For strictness, if $\mu\leq\mu_0$, then $I=\mathbb E(\mu_Y-\mu_0)_+>0$; if $\mu\geq\mu_0$, then $I=\mathbb E(\mu_0-\mu_Y)_+>0$. The last assertion follows from (1).
\end{proof}

\begin{example}[An exact prior]
Let $\Theta$ equal $0.9$ or $0.1$, each with probability one half, and let $\mu_0=0.6$. Then $\mu=0.5$. Bayes' rule yields posterior means $0.82$ after success and $0.18$ after failure, each outcome having probability one half. Consequently
\[
 \mathbb E\max\{0.6,\mu_Y\}=0.71,
 \qquad G=-0.10+0.11\beta.
\]
A planner with $\beta>10/11$ chooses the direction with the lower contemporaneous success probability. At $\beta=1$ the gain is $0.01$. These are exact arithmetic consequences of the specified prior, not empirical probabilities.
\end{example}

\paragraph{A contractor, capture, and a feasible grant.}
A one-period contractor conducts date-zero research and receives private success rewards $v_0,v_E\geq0$. Common real effort cost is $k\geq0$. Assume $v_0\mu_0-k\geq0$ and
\[
 d=v_0\mu_0-v_E\mu>0.
\]
Thus absent policy the contractor strictly chooses the safe project and participates. A grant $t$ is paid for the verified experimental direction and effort. A budget requires $0\leq t\leq B$. Grant financing has additional real cost $\lambda t$, $\lambda\geq0$; the grant itself cancels between taxpayer and domestic contractor. All success rewards are transfers already included in social output accounting.

An institutional failure occurs with exogenous probability $\rho\in[0,1)$. In that event verification is corrupted: the grant is disbursed, the safe project is conducted, and an additional real loss $L_{\mathrm{cap}}\geq0$ occurs. Otherwise verification is accurate. Institutional states are resolved before project choice, and this binary capture technology is common knowledge. It is an assumption about implementation, not a derived lobbying equilibrium.

\begin{proposition}[A budget and experimentation test]
Fix $\varepsilon\in(0,1]$. Suppose accurate-verification ties favor experimentation. A grant implementing experimental research whenever verification is accurate exists within the budget exactly when $B\geq d$. The smallest such grant is $d$. At that grant the unconditional experimentation probability is $1-\rho$ and the social gain is
\[
 (1-\rho)G-\rho L_{\mathrm{cap}}-\lambda d.
 \tag{2}
\]
Thus a weak implementation meeting the experimentation requirement and strictly improving welfare exists whenever
\[
 B\geq d,\qquad 1-\rho\geq\varepsilon,
 \qquad (1-\rho)G>\rho L_{\mathrm{cap}}+\lambda d.
\]
For strict contractor incentives use any $t>d$ satisfying the budget and replace $d$ by $t$ in (2).
\end{proposition}
\begin{proof}
Under accurate verification the experimental payoff minus the safe payoff is $t-d$. This proves the threshold, including the role of tie breaking and the strict variant. The safe payoff is nonnegative, so participation follows. Accurate verification produces social improvement $G$; capture reproduces the baseline project and subtracts $L_{\mathrm{cap}}$. The grant is paid in both states, generating real financing cost $\lambda t$. Taking expectations establishes (2). Choosing the smallest implementing grant weakly minimizes this cost and does not change the induced project or its information.
\end{proof}

\paragraph{Attempted extension and residual gap.}
The information term in (1) is a value-of-information calculation; no originality claim is made for Jensen's inequality. Its consequence for the policy target is sharper than a ranking by expected current output: even perfect observation of current success probabilities omits a potentially decisive public learning benefit. Equation (2) also shows why identifying that benefit alone is insufficient when verification can be captured.

With privately observed quality $q$, grant acceptance can select different priors and the public posterior need not equal the expression used here. A menu must then satisfy incentive constraints across types as well as directions and effort. Multiperiod experimentation changes the information state and may create endogenous abandonment or excessive continuation. Those are the unresolved parts of the catalogue problem. The explicit benchmark solves a grant threshold and an experimentation tradeoff under known priors and a specified capture process; it does not deliver a generally optimal innovation policy.

\begin{thebibliography}{9}
\bibitem{akcigit} U. Akcigit (2026), ``Creative destruction in economic growth,'' \emph{Scandinavian Journal of Economics}, Section 7, question 7. \href{https://onlinelibrary.wiley.com/doi/10.1111/sjoe.70037}{Publisher article}.
\end{thebibliography}

\end{document}
